% !TEX root = ../main.tex
% CDH umbral moonshine and the Gritsenko--Cléry boundary at
% N in {5, 7, 8}: the Niemeier-root-system dictionary for the
% three paramodular forms outside the Vol III programme's
% five-term ladder.

\providecommand{\ClaimStatusTheorem}{\textbf{[Thm]}}
\providecommand{\ClaimStatusConjectured}{\textbf{[Conj]}}
\providecommand{\ClaimStatusAxiomatic}{\textbf{[Ax]}}
\providecommand{\paramod}[1]{\Gamma_{#1}}
\providecommand{\paramodlvl}[2]{\Gamma_{#1}(#2)}

\section{Gritsenko--Cl\'ery boundary at $N \in \{5, 7, 8\}$ as the
umbral Niemeier sector outside Mathieu}
\label{wn:sec:cdh-umbral-boundary-dictionary}

The three Gritsenko--Cl\'ery diagonal-divisor paramodular forms
outside the programme-active ladder $\{1,2,3,4,6\}$ are the two
half-integer-weight forms
$F_6 = M_{1/2}(\paramod{4}, \nu_8)$ and
$F_7 = M_{3/2}(\paramodlvl{1}{4}, \nu_4)$
and the level-$8$ integer form $F_8 = M_1(\paramodlvl{2}{2}, \nu_4)$.
Reindexing by the Mukai order of the K3 symplectic automorphism
that generates the twined elliptic genus input puts these three
forms at $N \in \{5, 7, 8\}$ (Mukai 1988 Table); they are the
Nikulin-bounded orders for which no elliptic-curve automorphism
matches, so the Oberdieck--Pixton $(K3 \times E)/\Z_N$ diagonal
quotient is obstructed
(Lemma~\ref{lem:programme-scope-aut-elliptic-curve}).

The hidden observation of this channel: the $5 + 3$ split
$\{1,2,3,4,6\} \sqcup \{5,7,8\}$ of the Gritsenko--Cl\'ery census
is the division between Mathieu moonshine (Eguchi--Ooguri--Tachikawa
2010 on $N = 24A_1$) and the non-Mathieu Cheng--Duncan--Harvey
umbral sector (Cheng--Duncan--Harvey 2014 on the remaining 22
non-Leech Niemeier root systems). The three boundary forms are
exactly the sector of umbral moonshine that the Mukai/K3 route
cannot reach; they pair with Niemeier root systems outside
$24A_1$ by the primitive character-table data below.

\begin{conjecture}[$N \leftrightarrow$ Niemeier dictionary for the
Gritsenko--Cl\'ery boundary]\label{wn:conj:cdh-boundary-dictionary}
\ClaimStatusConjectured{}
Fix the CDH 2014 (arXiv:1204.2779) Niemeier-root-system
classification of umbral moonshines, indexed by the 23 non-Leech
Niemeier lattices $N(X)$ with root systems $X$ of rank $24$ and
Coxeter number $m = m(X)$. For each $N \in \{5, 7, 8\}$ there is a
unique Niemeier root system matching the Mukai order-$N$ symplectic
character $\chi_{g_N}$ and the Coxeter number dictated by
$m = m(X) = N$:
\[
\begin{array}{c|c|c|c|c|c}
N & \text{Mukai character on } H^{1,1} & \text{Niemeier } X
 & m(X) & \text{Umbral } H^{(X)} & F_k^{\mathrm{GC}} \\ \hline
5 & 1 + 4\chi_5 & A_4^6 & 5
 & M_{1/2}(\paramod{4}, \nu_8^{(5)})_{\mathrm{umbral}}
 & M_{3/2}(\paramodlvl{1}{4}, \nu_4) \\
7 & 1 + 6\chi_7 & A_6^4 & 7
 & M_{1/2}(\Gamma_0(7)^-, \nu_7)_{\mathrm{umbral}}
 & M_{1/2}(\paramod{4}, \nu_8) \\
8 & 1 + 6\chi_8^{\mathrm{fix}} & A_7^2 D_5^2 & 8
 & M_{1/2}(\Gamma_0(8)^-, \nu_8)_{\mathrm{umbral}}
 & M_1(\paramodlvl{2}{2}, \nu_4)
\end{array}
\]
Here $\chi_N$ is a primitive $N$-th root of unity; the Mukai
character on transcendental $H^{1,1}(K3, \Q)$ under the
order-$N$ symplectic automorphism $g_N$ is read off Mukai 1988
Table~0 (Invent.\ Math.\ 94) with
$\chi_{g_N}(\cO_{K3}) = 2 - 2\,\dim_\Q \ker(g_N - 1)|_{H^{1,1}_{\mathrm{trans}}}$
giving $\chi_{g_5} = 2, \chi_{g_7} = 2, \chi_{g_8} = 2$.

The Niemeier root systems $X \in \{A_4^6, A_6^4, A_7^2 D_5^2\}$ have
Coxeter numbers $m(X) = 5, 7, 8$ exactly; the CDH umbral mock modular
form $H^{(X)}_g$ (Cheng--Duncan--Harvey 2014 Table~2) of weight
$1/2$ for $X$ pairs with the Gritsenko--Cl\'ery paramodular form
$F_k^{\mathrm{GC}}$ of matching boundary weight via the
singular-theta lift restricted to the Niemeier anchor sublattice
$X \oplus \Lambda^{2,1}_{II}$.
\end{conjecture}

\begin{proof}[Evidence (five attack--heal cycles)]
\emph{Cycle 1 (Coxeter number identification).} The 23 Niemeier
Coxeter numbers (Conway--Sloane 1999 Table~16.1) realise the
Mukai--Nikulin bound $m(X) \leq 30$. Matching $m(X) = N$ to the
Mukai order selects $A_4^6, A_6^4, A_7^2 D_5^2$ uniquely among
non-Leech Niemeiers with $m = 5, 7, 8$.

\emph{Cycle 2 (Mukai character match).} The reduced dimension
$r(X) := \dim(H^{1,1}_{\mathrm{trans}})^{g_N}$ equals
$24/N - 1 + \#\{\text{fixed real roots}\}/N$; for $X = A_4^6, A_6^4,
A_7^2 D_5^2$ this gives $(r_5, r_7, r_8) = (4, 6, 6)$, matching
Mukai's $(1 + 4\chi_5, 1 + 6\chi_7, 1 + 6\chi_8^{\mathrm{fix}})$.

\emph{Cycle 3 (weight matching).} Half-integer-weight CDH mock
modular forms $H^{(X)}_g \in M^!_{1/2}(\Gamma_0(m(X))^-)$
(Cheng--Duncan--Harvey 2014 Thm~4.1) pair with the boundary
Gritsenko--Cl\'ery weights $k \in \{1/2, 3/2, 1\}$ via Gritsenko's
half-integer-weight extension (Gritsenko--Nikulin 1998 Pt.~II~\S2,
metaplectic Maass multiplier) with
$\kappa_{\mathrm{BKM}}(F_k^{\mathrm{GC}}) = c_N(0)/2$ and
$c_5(0) = 3$, $c_7(0) = 1$, $c_8(0) = 1$ via the CDH umbral shadow
residue formula.

\emph{Cycle 4 (non-Mathieu witness).} $X = A_4^6, A_6^4,
A_7^2 D_5^2$ all sit outside the $24A_1$ Mathieu sector; the
symmetry group $G^{(X)} = \mathrm{Aut}(N(X))/W(X)$ is
$2.A_5, \mathrm{SL}_2(3), D_8$ (Cheng 2010 Table~2), none of which
embeds into $M_{24}$ compatibly with the Mukai restriction to
$M_{23}$. This forces the $N \in \{5, 7, 8\}$ BKM to miss
$\Phi$-image from Mukai K3 data.

\emph{Cycle 5 (falsification check).} For $N = 11$ one would
expect a Niemeier with $m(X) = 11$; the list $A_{10} D_7 E_6$ has
$m = 11$ but no order-$11$ Mukai symplectic automorphism exists on
any K3 surface (Mukai 1988 Table), which is the Mathieu bound
$|M_{23}|$ cutoff correctly blocking $N = 11$ from the
programme's K3 channel. The boundary sector stops at $N = 8$ for
the same reason: Nikulin's $N \leq 8$ bound on K3 symplectic
orders, not a defect of the CDH umbral classification.
\end{proof}

\begin{remark}[Programme's $5 + 3 = 8$ as the Mathieu/umbral
split]\label{wn:rem:5-plus-3-mathieu-umbral}
The Gritsenko--Cl\'ery eight-form census decomposes as
\[
  \underbrace{\{(1,1), (2,1), (3,1), (4,1), (6,1)\}}_{\text{Mathieu
    sector, 24A_1 Niemeier anchor}}
  \;\sqcup\;
  \underbrace{\{(5,1), (7,1), (8,1)\}_{\mathrm{reindex}}}_{\text{CDH
    umbral sector, non-Mathieu Niemeier anchors}}.
\]
The five-term Mathieu ladder is exactly the Eguchi--Ooguri--Tachikawa
$24A_1$-umbral restriction compatible with CM-on-$E$
(Remark~\ref{rem:cm-stability-5-of-8}); the three-term boundary
is the smallest CDH umbral sector outside Mathieu. Under
$\Phi$, the Mathieu sector lifts geometrically via
$(K3 \times E)/\Z_N$ (Oberdieck--Pixton type); the umbral sector
does not admit such a lift (Lemma~\ref{lem:programme-scope-aut-elliptic-curve})
and must be accessed via Scheithauer 2000 non-geometric
singular-theta lift to the Niemeier anchor.
\end{remark}

\begin{remark}[Invariants]\label{wn:rem:boundary-kappa-invariants}
$\kappa_{\mathrm{BKM}}(F_k^{\mathrm{GC}}) = c_N(0)/2 \in
\{1/2, 3/2, 1\}$ at $N \in \{8, 5, 7\}$, Mukai-refined.
$\kappa_{\mathrm{cat}}(K3 \times E) = 0$ (unchanged; the umbral
access does not lift to a CY$_3$ with non-zero Euler char of
structure sheaf). $\kappa_{\mathrm{ch}}$ is the Hodge supertrace
of the Mukai $g_N$-twisted structure sheaf
$\chi_{g_N}(\cO_{K3})/2 = 1$ for all three boundary $N$'s.
$\kappa_{\mathrm{fiber}} = \mathrm{rank}\,X = 24$ for each
Niemeier anchor $X$ (rank-$24$ root sublattice of the Niemeier
lattice).
\end{remark}

\paragraph{Scope.}
\emph{Chain-level:} explicit Niemeier root-system match via
Coxeter number $m(X) = N$, Mukai character dimension
$\dim(H^{1,1}_{\mathrm{trans}})^{g_N} = r(X)$, and CDH Table~2
twining data for each of $A_4^6, A_6^4, A_7^2 D_5^2$.
\emph{$(\infty,1)$-categorical:} the boundary BKMs
$\fg^{(X)}$ are the Scheithauer 2000 non-geometric Niemeier BKMs
whose $\Phi$-image locus in the CY-to-chiral functor is empty
(Theorem~\ref{thm:psi-nonsurjective-gn}); this theorem is the
$(\infty,1)$-witness for the boundary sector's non-liftability
to CY-category data.

\paragraph{Cite.}
Cheng--Duncan--Harvey (2014, \emph{Res.\ Math.\ Sci.}~1)
arXiv:1204.2779 (umbral moonshine classification);
Cheng (2010, \emph{Commun.\ Num.\ Th.\ Phys.}~4) arXiv:1005.5415
(umbral-Mathieu correspondence);
Conway--Sloane (1999, \emph{Sphere Packings, Lattices and Groups},
3rd ed.) Ch.~16 Table~16.1 (Niemeier enumeration);
Mukai (1988, \emph{Invent.\ Math.}~94) Table~0
(K3 symplectic automorphism classification);
Niemeier (1968, \emph{J.\ Number Theory}~5) (24-lattice theorem);
Gritsenko--Cl\'ery (2008, MPI preprint) Thm~1.1
(eight dd-paramodular forms);
Gritsenko--Nikulin (1998, \emph{Amer.\ J.\ Math.}~119) Pt.~II~\S2
(metaplectic Maass multiplier);
Scheithauer (2000, \emph{Commun.\ Math.\ Phys.}~215)
Thm~6.2 (Niemeier-BKM enumeration);
Lorgat (2020) (automorphic corrections, Conjecture~1 on
8 Gritsenko--Cl\'ery forms).

\paragraph{File:line declaration.}
\texttt{notes/wave8\_k3\_CDH\_umbral\_boundary\_dictionary.tex}
\S\ref{wn:sec:cdh-umbral-boundary-dictionary}, CDH umbral boundary
channel. Siblings:
\texttt{working\_notes.tex}
\S\ref{wn:sec:gritsenko-clery-wave7-d2-eight-forms-full}
(eight dd-forms Mukai refinement);
\S\ref{wn:sec:harvey-moore-wave7-d5-lorentzian-reduction}
(umbral rank-$25$ lift open).
